\scanpage{119}
% archive foliation: 118
satisfaisant à la condition que son carré $\kappa^{\unsure{\otimes2|\otimes}}$ s'\deleted{\unsure{se}}identifie, via les isom. canoniques
\[
\tag{77}
\underline{t}_E^{\unsure{\otimes2|\otimes}}\simeq T_{\Sigma_J,t},
\qquad
\mathcal O_t(1)^{\otimes2}=\mathcal O_t(2)
\underset{\text{via }\alpha}{\simeq}
\mathcal O_t(2)(\underline{\omega})\simeq T_{\Sigma_J,t}
\]
à l'isom identique de $T_{\Sigma_J,t}$ --- ce qui, pour $\alpha$ fixé (deux choix), détermine $\kappa$ modulo signe. Pour $\alpha$ non fixé, $\kappa$ est déterminé modulo multiplication par $\mu_4(k)$.

Rappelons-nous de la suite exacte canonique

(78)
\diagram{\begin{tikzcd}[ampersand replacement=\&,column sep=small]
0 \arrow[r] \& F_t \arrow[r] \arrow[d,"\wr" description,phantom] \&
V_J(k) \arrow[r] \arrow[d,"\wr" description,phantom] \&
\mathcal O_t(1) \arrow[r] \& 0\\
\& \mathcal O_t(-1)(\underline{\omega}) \&
\Gamma(\Sigma_J,\underline{\mathcal O}_{\Sigma_J}(1)) \&\&
\end{tikzcd}}

que \deleted{\unsure{pour faire}} je compare à

(79)
\diagram{\begin{tikzcd}[ampersand replacement=\&,column sep=small]
0 \arrow[r] \& k.1_{\Sigma_J} \arrow[r] \&
\overbrace{\Gamma(\Sigma_J,\underbrace{\underline{\mathcal O}_{\Sigma_J}\{t\}}_{\substack{\simeq\underline{\mathcal O}_{\Sigma_J}(1)\\\text{non canoniquement}}})}^{\simeq F_2}
\arrow[r] \& T_{\Sigma_J,t}\arrow[r] \& 0
\end{tikzcd}}

On constate que la deuxième suite exacte est canoniquement isom. à celle déduite de la première, en la tensorisant par
\[
\tag{80}
\check F_t\simeq\mathcal O_t(1)(\underline{\omega})
\]
(on a un iso valable chaque fois qu'on a un Module loc libre de rang $2$, $\underline V$, \deleted{\unsure{dont}} et une section $t$ de $\mathbb P(\underline V)$ \ldots) :

\deleted{(79)} (81)
\diagram{\begin{tikzcd}[ampersand replacement=\&,column sep=small]
0 \arrow[r,dash] \&
F_t\otimes\check F_t \arrow[r] \arrow[d,"\wr" description] \&
V_J(k)\otimes\check F_t \arrow[r] \arrow[d,"\wr" description] \&
\overbrace{\mathcal O_t(1)\otimes\check F_t}^{\underline{\mathcal O}_t(2)(\underline{\omega})} \arrow[r] \arrow[d,"\wr" description] \& 0\\
0 \arrow[r] \& k.1_{\Sigma_J}\arrow[r] \&
\underbrace{\Gamma(\Sigma_J,\underline{\mathcal O}_{\Sigma_J}\{t\})}_{F_2}
\arrow[r] \& T_{\Sigma_J,t}\arrow[r] \& 0 .
\end{tikzcd}}
